\documentclass[10pt,letterpaper]{article}
\usepackage[margin=0.88in]{geometry}
\usepackage[T1]{fontenc}
\usepackage{lmodern}
\usepackage{amsmath,amssymb,amsthm,booktabs}
\usepackage[colorlinks=true,linkcolor=blue,citecolor=blue,urlcolor=blue]{hyperref}
\newtheorem{theorem}{Theorem}
\newtheorem{lemma}{Lemma}
\newtheorem{corollary}{Corollary}
\DeclareMathOperator{\tr}{tr}
\DeclareMathOperator{\diag}{diag}
\newcommand{\E}{\mathbb E}
\newcommand{\RP}{\mathrm{RP}}
\newcommand{\DPP}{\mathrm{DPP}}
\newcommand{\TV}{d_{\mathrm{TV}}}
\newcommand{\R}{\mathbb R}
\newcommand{\Z}{\mathbb Z}
\allowdisplaybreaks[1]
\raggedbottom
\title{Sharpness of RPCholesky error doubling\\with an explicit integer family}
\author{\small AI-assisted research project directed by Jitendra Prajapati}
\date{18 October 2026 --- public research draft, version 2}

\begin{document}
\maketitle
\begin{abstract}
Chen, Epperly, Tropp, and Webber proved that the expected trace residual
after exactly $k$ randomly pivoted Cholesky pivots is at most $2^k$ times
the optimal rank-$k$ trace error. We give an explicit positive-definite
family in dimension $k+1$ approaching this factor. A single rational
parameter, explicit error bounds, and an integer rescaling yield finite
witnesses with polynomial encoding length. On the same family, the
normalized expected residual of size-$k$ determinant sampling tends to
$k+1$. The construction concerns exactly $k$ pivots; one further pivot
exhausts the family, so no oversampling conjecture is refuted. The proposed
matching lower bound has undergone separately tasked AI checks, but human
mathematical review, novelty, and publication priority remain unconfirmed.
\end{abstract}

\section{Statement, algorithm, and attribution}

For a real symmetric positive-definite $d\times d$ matrix $A$, ordinary
randomly pivoted Cholesky (RPCholesky) starts with $R_0=A$. Given $R_j$,
it chooses $i$ with probability $(R_j)_{ii}/\tr R_j$ and sets
\[
 R_{j+1}=R_j-\frac{R_je_ie_i^{\mathsf T}R_j}{(R_j)_{ii}}.
\]
An exhausted residual stays zero. Let
$\tau_k(A)=\sum_{i>k}\lambda_i(A)$, with eigenvalues in decreasing order.
Thus $\tau_k$ is the optimal rank-$k$ trace error. All statements below
use exact arithmetic and this prescribed pivot law. They concern an
expectation over pivot choices, not a lower bound on every sampled path.

\begin{theorem}[Sharpness of the published upper bound]\label{thm:sharp}
For every integer $k\geq1$,
\begin{equation}\label{eq:supremum}
 \sup_{d>k,\ A\succ0}\frac{\E\tr R_k(A)}{\tau_k(A)}=2^k.
\end{equation}
The supremum is unchanged when $A$ is restricted to real symmetric
positive-definite integer matrices of dimension $k+1$.
\end{theorem}

The upper bound is the existing result of Chen, Epperly, Tropp, and Webber
\cite[Lemma~5.5]{CETW}; it is not a contribution of this note.
The proposed contribution is a matching explicit lower construction and
its quantitative analysis. RPCholesky, fixed-size determinant sampling,
Schur-complement identities, Cauchy--Binet, and cofactor identities are
established ingredients. Perturbed-simplex lower bounds, including the
$k=1$ limiting factor $2$, are established in \cite[Proposition~3.3]{DRVW}.
Constructions combining orthogonal simplices with separated scales are
also established \cite[Section~3]{DKM}. The final inverse-diagonal and
determinant-sampling identities used here are standard and are recorded,
for example, in \cite[Section~4.2]{Simons}. No novelty is claimed for
these ingredients, the weighted contrast basis, or the general idea of
scale separation. The adaptive survivor-law analysis below concerns the
specific ordinary RPCholesky sampling law.
Theorem~\ref{thm:sharp} is independent of any interpretation of a
conjecture; the carefully qualified thesis implication appears in
Section~\ref{sec:scope}.

The proof will establish a rational one-parameter family with a uniform
error estimate, then clear its denominators. This avoids a choice of
successively smaller parameters without an explicit dependence between
them. In particular, rare pivot paths are controlled after multiplication
by their potentially large terminal amplification.

\section{The explicit family and quantitative assertion}

Fix $n=k+1\geq2$ and $0<t\leq1/2$. Define
\begin{equation}\label{eq:parameters}
 w_i=t^{i-1},\qquad S_m=\sum_{i=1}^m w_i^2,\qquad
 Z=S_n,\qquad p_i=\frac{w_i^2}{Z},\qquad
 h=t^{4n},\qquad \delta=t^{4n(n-1)}.
\end{equation}
For $1\leq m<n$, let the weighted contrast $v_m\in\R^n$ have entries
\begin{equation}\label{eq:vectors}
 (v_m)_i=\begin{cases}
 w_iw_{m+1},&1\leq i\leq m,\\
 -S_m,&i=m+1,\\
 0,&i>m+1.
 \end{cases}
 \qquad
 A_0=\sum_{m=1}^{n-1}h^{m-1}v_mv_m^{\mathsf T},\qquad
 A=A_0+\delta I_n.
\end{equation}
The arguments in Sections~\ref{sec:inputs}--\ref{sec:noise} first allow
arbitrary $0<h\leq1$ and $\delta>0$; the values in
\eqref{eq:parameters} are substituted only at the end.

\begin{theorem}[Uniform finite-parameter estimate]\label{thm:quantitative}
The matrix $A$ in \eqref{eq:parameters}--\eqref{eq:vectors} is positive
definite, rational for rational $t$, and has $\tau_{n-1}(A)=\delta$.
Moreover,
\begin{align}
 \left|\frac{\E\tr R_{n-1}(A)}{\delta}-2^{n-1}\right|
 &\leq n2^{n-1}t^2+
 \bigl[2^n n(n-1)+2(1+(n-1)^2)\bigr]t^4\label{eq:finebound}\\
 &\leq 2^{n+1}n^2t^2.\label{eq:simplebound}
\end{align}
Consequently, for $\eta>0$, any integer
\begin{equation}\label{eq:qchoice}
 q\geq\max\{2,2n/\sqrt\eta\},\qquad t=1/q,
\end{equation}
gives a finite rational matrix satisfying
\begin{equation}\label{eq:relative}
 (1-\eta)2^{n-1}\leq
 \frac{\E\tr R_{n-1}(A)}{\tau_{n-1}(A)}\leq2^{n-1}.
\end{equation}
The last upper bound invokes \cite[Lemma~5.5]{CETW}.
\end{theorem}

The constants and exponents are convenient safe choices; optimality of
these choices is not asserted. For absolute error $\epsilon$ in the
normalized expectation, it suffices to replace \eqref{eq:qchoice} by
$q\geq\max\{2,\sqrt{2^{n+1}n^2/\epsilon}\}$.

\section{Uniform spectral and principal-minor inputs}\label{sec:inputs}

For $m<\ell$, the restriction of $v_\ell$ to the support of $v_m$ is
$w_{\ell+1}w$. Since $v_m^{\mathsf T}w=0$, the vectors $v_m$ are
nonzero and mutually orthogonal. Directly,
$\|v_m\|^2=S_mS_{m+1}$. Hence $A_0$ has kernel $\operatorname{span}(w)$,
and its nonzero eigenvalues are
\begin{equation}\label{eq:eigenvalues}
 \lambda_m=h^{m-1}S_mS_{m+1},\qquad
 \lambda_*:=\min_m\lambda_m\geq h^{n-2}.
\end{equation}
No ordering of the individual $\lambda_m$ is required. Adding $\delta I$
makes the smallest eigenvalue exactly $\delta$, proving the spectral-tail
assertion of Theorem~\ref{thm:quantitative}.

Write $u=w/\sqrt Z$ and $p_{\min}=p_n$. The semidefinite inequality
$A_0\succeq\lambda_*(I-uu^{\mathsf T})$ gives, for every nonempty
proper coordinate set $T$,
\begin{equation}\label{eq:principal-lower}
 A_0[T,T]\succeq h^{n-2}p_{\min}I_T=:\mu I_T.
\end{equation}
Indeed the least eigenvalue of $I_T-u_Tu_T^{\mathsf T}$ is
$1-\|u_T\|^2=\sum_{j\notin T}p_j\geq p_{\min}$.
Thus every proper principal submatrix is positive definite. Also
\begin{equation}\label{eq:pmin}
 1\leq Z\leq\frac43,\qquad
 p_{\min}^{-1}\leq2t^{-2(n-1)},\qquad
 p_{\min}^{-2}\leq2t^{-4(n-1)}.
\end{equation}
The last constant follows from $Z^2\leq16/9<2$, independently of the
preceding, looser factor-$2$ bound.

Let $V$ have rows $v_1,\ldots,v_{n-1}$. For $r\leq n-1$, restrict its
first $r$ rows to the first $r+1$ coordinates, obtaining $H_r$.
Its kernel is spanned by $(w_1,\ldots,w_{r+1})$, so its signed maximal
cofactors are proportional to these entries. Cauchy--Binet determines the
constant:
\begin{equation}\label{eq:cofactors}
 \det(H_r\text{ with column }j\text{ omitted})^2
   =w_j^2\left(\prod_{\ell=1}^r S_\ell\right)^2.
\end{equation}
For completeness, the sum of the squared cofactors is
$\det(H_rH_r^{\mathsf T})=\prod_{\ell=1}^r S_\ell S_{\ell+1}$.
Divide by $\sum_{j=1}^{r+1}w_j^2=S_{r+1}$ and use $S_1=1$ to obtain
\eqref{eq:cofactors}.

An additional coefficient-sum bound will avoid any hidden combinatorial
factor. Put $B=\sum_{m=1}^{n-1}v_mv_m^{\mathsf T}$. Its nonzero
eigenvalues are $S_mS_{m+1}$, so $\|B\|\leq16/9<2$. For every
size-$r$ coordinate set $T$,
\begin{equation}\label{eq:minor-sum}
 \sum_{\substack{I\subseteq\{1,\ldots,n-1\}\\|I|=r}}
       \det(V[I,T])^2
   =\det B[T,T]\leq2^r.
\end{equation}
This bounds the sum of all squared-minor coefficients, not just each
coefficient separately.

\section{The hierarchical law and its weighted error}\label{sec:hierarchy}

After selecting a coordinate set $S$, the residual diagonal at unused
$i$ equals $\det A_0[S\cup\{i\},S\cup\{i\}]/\det A_0[S,S]$.
Consequently the next-pivot law on $A_0$ is
\begin{equation}\label{eq:transition}
 P_0(i\mid S)=\frac{\det A_0[S\cup\{i\},S\cup\{i\}]}
 {\sum_{j\notin S}\det A_0[S\cup\{j\},S\cup\{j\}]}.
\end{equation}
This is a transition identity; it does not identify the final RP subset
law with determinant sampling. All blocks used have size at most $n-1$.

Call a size-$m$ selected set \emph{legal} if
$S=\{1,\ldots,m+1\}\setminus\{J\}$ for some $J\leq m+1$.
The empty set is legal with representative $J=1$. At a legal state with
$m\leq n-2$, put $r=m+1$ and $e=r(r-1)/2$. Cauchy--Binet gives
\begin{equation}\label{eq:cb}
 \det A_0[T,T]=\sum_{|I|=r}
 h^{\sum_{a\in I}(a-1)}\det(V[I,T])^2,
 \qquad |T|=r.
\end{equation}
The unique least exponent is $e$, from $I=\{1,\ldots,r\}$;
every other exponent is at least $e+1$. For $T=S\cup\{i\}$, the
leading coefficient is nonzero only for $i=J$ or $i=m+2$.
By \eqref{eq:cofactors}, these coefficients are
\begin{equation}\label{eq:leading}
 c_J=w_{m+2}^2\left(\prod_{\ell=1}^r S_\ell\right)^2,
 \qquad
 c_{m+2}=w_J^2\left(\prod_{\ell=1}^r S_\ell\right)^2.
\end{equation}
For other unused $i$, set $c_i=0$. Write each determinant weight as
$h^e c_i+E_i$. Nonnegative coefficients, $h\leq1$, and
\eqref{eq:minor-sum} imply $0\leq E_i\leq2^rh^{e+1}$.
Thus the total leading and remainder weights obey
\begin{equation}\label{eq:leading-remainder}
 L:=h^e\sum_i c_i\geq h^e t^{2(n-1)},\qquad
 E:=\sum_iE_i\leq n2^{n-1}h^{e+1}.
\end{equation}
The lower bound is deliberately loose: $J\leq n-1$ and each $S_\ell\geq1$.

Let $Q$ normalize the two leading coefficients. The actual transition is
exactly the mixture $(L/(L+E))Q+(E/(L+E))U$, where $U$ normalizes the
remainders when $E>0$. Its total variation distance from $Q$ is at most
\begin{equation}\label{eq:alpha}
 \alpha:=n2^{n-1}h\,t^{-2(n-1)}.
\end{equation}
We use $\TV$ as the maximal difference of event probabilities, equal to
the smallest possible disagreement probability in a coupling. Couple
the RP path and the $Q$ path while they agree. The latter always stays
legal; a union bound over $n-1$ transitions bounds the probability of any
disagreement by $(n-1)\alpha$. This includes all illegal RP paths, and
requires no estimate after the first disagreement.

Let $J$ now denote the final unselected coordinate and put $f(J)=1/p_J$.
Because $0<f\leq1/p_{\min}$, \eqref{eq:pmin} gives the crucial weighted
bound
\begin{equation}\label{eq:hierarchy-error}
 \left|\E_{P_0}f(J)-\E_Q f(J)\right|
 \leq2^n n(n-1)h\,t^{-4(n-1)}.
\end{equation}
In particular, convergence of transition probabilities alone has not
been used to justify convergence of this unbounded-as-$t\to0$ payoff.

The law $Q$ is an explicit tournament. Starting with $J_1=1$, introduce
$j=2,\ldots,n$; if the old representative is $i<j$, it survives with
probability $p_i/(p_i+p_j)$, and the new representative is $j$ otherwise.
The opposite-column coefficients in \eqref{eq:leading} give exactly
these survival probabilities. If $F_j=\E[1/p_{J_j}]$, then
\[
 \E[1/p_{J_j}\mid J_{j-1}=i]=\frac2{p_i+p_j},
 \qquad F_1=Z.
\]
Since $p_j/p_i=t^{2(j-i)}\leq t^2$,
\begin{equation}\label{eq:tournament}
 Z\left(\frac2{1+t^2}\right)^{n-1}\leq F_n\leq Z2^{n-1},
 \qquad |F_n-2^{n-1}|\leq n2^{n-1}t^2.
\end{equation}
For the last inequality use $(1+x)^{-(n-1)}\geq1-(n-1)x$ and
$(1-t^2)^{-1}\leq1+\frac43t^2$ for $t\leq1/2$.

\section{Uniform noise control and completion of the bound}\label{sec:noise}

Let $P_\delta$ be the RP path law for $A_0+\delta I$. For a proper set
$T$ of size $r$, \eqref{eq:principal-lower} implies, with $x=\delta/\mu$,
\[
 1\leq\frac{\det(A_0[T,T]+\delta I_T)}{\det A_0[T,T]}
 \leq(1+x)^r.
\]
Every transition weight in \eqref{eq:transition} is therefore tilted by
a factor between $1$ and $\beta=(1+x)^{n-1}$. It follows pointwise that
$P_\delta(i\mid S)\geq P_0(i\mid S)/\beta$. The transition-law overlap
is at least $\beta^{-1}$, so its total variation distance is at most
$1-\beta^{-1}\leq(n-1)x$. This last inequality holds for every $x\geq0$.
Coupling all $n-1$ transitions, now at arbitrary selected states, gives
\begin{equation}\label{eq:noise-law}
 \TV(P_\delta,P_0)\leq(n-1)^2\frac{\delta}{\mu}.
\end{equation}

After $n-1$ pivots, the scalar Schur-complement identity gives
$\tr R_{n-1}(A)=1/(A^{-1})_{JJ}$. Spectral decomposition yields
\[
 (A^{-1})_{jj}=p_j/\delta+b_j,
 \qquad 0\leq b_j\leq\lambda_*^{-1}\leq h^{-(n-2)}.
\]
Thus the normalized endpoint residual is
$g_\delta(j)=1/(p_j+\delta b_j)$, and
\begin{equation}\label{eq:endpoint}
 0\leq f(j)-g_\delta(j)
 \leq\frac{\delta h^{-(n-2)}}{p_{\min}^2}.
\end{equation}
Comparing $\E_{P_\delta}g_\delta$ first with $\E_{P_\delta}f$ and
then with $\E_{P_0}f$, equations \eqref{eq:noise-law}--\eqref{eq:endpoint}
and $\mu=h^{n-2}p_{\min}$ show
\begin{align}
 \left|\E_{P_\delta}g_\delta(J)-\E_{P_0}f(J)\right|
 &\leq[1+(n-1)^2]\frac{\delta h^{-(n-2)}}{p_{\min}^2}\notag\\
 &\leq2[1+(n-1)^2]\delta h^{-(n-2)}t^{-4(n-1)}.
 \label{eq:noise-error}
\end{align}
The comparison payoff $f$ under $P_0$ is a limiting amplification weight;
it is not the residual trace of the rank-$(n-1)$ matrix $A_0$, which
vanishes after $n-1$ pivots. This distinction accounts for both endpoint
and path-law perturbations.

\begin{proof}[Completion of Theorem~\ref{thm:quantitative}]
Substitute $h=t^{4n}$ and $\delta=t^{4n(n-1)}$. The two dangerous
weighted factors become
\begin{equation}\label{eq:exponents}
 h t^{-4(n-1)}=t^4,\qquad
 \delta h^{-(n-2)}t^{-4(n-1)}=t^4.
\end{equation}
The triangle inequality applied to \eqref{eq:hierarchy-error},
\eqref{eq:tournament}, and \eqref{eq:noise-error} proves
\eqref{eq:finebound}. Since $t^4\leq t^2$ and
\[
 n2^{n-1}+2^n n(n-1)\leq2^n n^2,
 \qquad 2[1+(n-1)^2]\leq2n^2\leq2^n n^2,
\]
we obtain \eqref{eq:simplebound}. Dividing by $2^{n-1}$ and taking
$t=1/q$ gives relative error at most $4n^2/q^2$, proving the tolerance
claim. The published upper bound supplies the upper side of
\eqref{eq:relative}.
\end{proof}

\section{Integer matrices and encoding length}\label{sec:integer}

The rational family admits a formula using only integer powers and sums.
Let $q\geq2$ be an integer, $t=1/q$, and $D=4n(n-1)$. Define
\begin{equation}\label{eq:integer-vectors}
 (z_m)_i=\begin{cases}
 q^{m-i},&1\leq i\leq m,\\
 -\displaystyle\sum_{\ell=0}^{m-1}q^{2m-1-2\ell},&i=m+1,\\
 0,&i>m+1.
 \end{cases}
\end{equation}
These are precisely $z_m=q^{2m-1}v_m$. Therefore
\begin{equation}\label{eq:integer-matrix}
 M_{n,q}:=I_n+\sum_{m=1}^{n-1}
 q^{D-4n(m-1)-4m+2}z_mz_m^{\mathsf T}=q^D A.
\end{equation}
The coefficient exponent decreases with $m$ and equals $6$ at $m=n-1$;
all coefficients are positive integers. The $z_m$ are orthogonal and
orthogonal to $(q^{n-1},q^{n-2},\ldots,q,1)$. Hence $M_{n,q}$ is
integer SPD, its least eigenvalue is $1$, and $\tau_{n-1}(M_{n,q})=1$.

Multiplication by a positive scalar leaves every RP pivot probability
unchanged and scales every residual and spectral tail by that scalar.
Theorem~\ref{thm:quantitative} therefore gives the following consequence.

\begin{corollary}[Finite integer witness]\label{cor:integer}
For $n\geq2$, $\eta>0$, and an integer $q$ satisfying
\eqref{eq:qchoice}, the explicit integer matrix \eqref{eq:integer-matrix}
satisfies
\[
 \tau_{n-1}(M_{n,q})=1,\qquad
 (1-\eta)2^{n-1}\leq\E\tr R_{n-1}(M_{n,q})\leq2^{n-1}.
\]
Each entry has $O(n^2\log q)$ bits, and the full dense matrix has
$O(n^4\log q)$ bits. With $0<\eta<1$ and
$q=\lceil2n/\sqrt\eta\rceil$, this is
$O(n^4(\log n+\log(1/\eta)))$ bits in total.
\end{corollary}

\begin{proof}
Only the size bound remains. Since $h\leq1$, $A_0\preceq B$, and
\eqref{eq:parameters} gives $\delta\leq1$. Thus
$\|A\|\leq16/9+1<3$. Every entry of $M_{n,q}=q^DA$ has absolute
value less than $3q^D$, proving the per-entry and dense encoding bounds.
The stated choice of $q$ has
$\log q=O(\log n+\log(1/\eta))$.
\end{proof}

This is a polynomial encoding-length statement, not a polynomial-time
claim for exact expectation evaluation. The small finite checks below
enumerate pivot paths. No numerical conditioning or floating-point
stability guarantee follows from a short exact encoding.

\begin{proof}[Proof of Theorem~\ref{thm:sharp}]
The published upper bound applies to every SPD input. For fixed $k$, put
$n=k+1$ in Corollary~\ref{cor:integer}, and let $\eta\downarrow0$.
Every chosen tolerance gives a finite integer SPD matrix, so the supremum
over that restricted class is at least $2^k$. Together with the upper
bound this proves both assertions. No singular limiting matrix is used
as an input, and no finite attainment of the supremum is asserted.
\end{proof}

\section{Same-size determinant sampling}

A size-$k$ determinant sample (a $k$-DPP with kernel $A$) selects $T$ with
probability $\det A[T,T]/e_k(A)$, where $e_k$ is the $k$th elementary
symmetric polynomial of the eigenvalues. Its residual is the same
selected-set Schur residual. For $k=n-1$, every set omits one $j$.
Multiplying its weight by its residual cancels its principal determinant,
giving the exact identity
\begin{equation}\label{eq:dpp-exact}
 \E\tr R_{n-1}^{\DPP}(A)
 =\frac{n\det A}{e_{n-1}(A)}=\frac n{\tr(A^{-1})}.
\end{equation}
By \eqref{eq:eigenvalues},
\[
 \frac{\E\tr R_{n-1}^{\DPP}(A)}{\delta}
 =\frac n{1+\delta\sum_{m=1}^{n-1}(\lambda_m+\delta)^{-1}},
 \qquad
 0\leq\delta\sum_m(\lambda_m+\delta)^{-1}
 \leq(n-1)t^{4n}.
\]
Consequently
\begin{equation}\label{eq:dpp-bounds}
 \frac n{1+(n-1)t^{4n}}
 \leq\frac{\E\tr R_{n-1}^{\DPP}(A)}{\delta}\leq n,
 \qquad
 \lim_{t\downarrow0}\frac{\E\tr R_{n-1}^{\RP}(A)}
 {\E\tr R_{n-1}^{\DPP}(A)}=\frac{2^{n-1}}n.
\end{equation}
For $0<\eta<1$, the integer tolerance choice gives a finite RP/DPP
expected-error ratio at least $(1-\eta)2^{n-1}/n$. Thus no universal
constant or polynomial in $k$ bounds the same-$k$ ratio for all SPD
matrices. The DPP value $k+1$ is a family limit, not an equality asserted
for every finite member. This compares trace errors at the same number
of selected coordinates, not runtimes or practical kernel performance.

\section{Illustrative exact finite checks}

The accompanying integer-family records contain ten identity cases
($n=2,\ldots,6$, $q=2,3$) and four complete exact expectation cases,
all recorded as passing the fixed protocol. The identity checks compare
\eqref{eq:integer-matrix} entry by entry with $q^DA$ and check its full
orthogonal eigenbasis. For the expectation cases, the smallest admissible
integer $q$ in \eqref{eq:qchoice} was chosen with $\eta=1/10$.
The recorded values are:
\begin{center}
\begin{tabular}{rrrrr}
\toprule
$k$ & $q$ & $\E\tr R_k(M_{k+1,q})$ & Ratio to $2^k$ & Max. entry bits\\
\midrule
1 & 13 & 1.9999999976 & 0.9999999988 & 30\\
2 & 19 & 3.9944903578 & 0.9986225894 & 102\\
3 & 26 & 7.9881918819 & 0.9985239852 & 226\\
4 & 32 & 15.9765967765 & 0.9985372985 & 401\\
\bottomrule
\end{tabular}
\end{center}
Here $\tau_k=1$ exactly. Displayed decimals are rounded; all pass/fail
comparisons use rational arithmetic. The direct recursion weights every
positive residual diagonal by its fraction of the trace, performs the
Schur update, and enumerates the remaining paths. A known $3\times3$
control gives $23/28$; an intentionally uniform-pivot mutation gives a
different result and is rejected. The saved finite expectations are
strictly below $2^k$.

These checks illustrate the construction and test the implementation.
They are not Monte Carlo estimates, benchmarks, or evidence by finite
extrapolation for arbitrary $n$. The theorem and its tolerance guarantee
rest on the proof above. A separate AI audit checked code, saved-record
consistency, and hashes without independently rerunning the expectation
recursion; it should not be described as an additional execution.

\section{Scope, thesis interpretation, and review status}\label{sec:scope}

The sharpness theorem rules out a universal polynomial $C(r)$ satisfying
\begin{equation}\label{eq:poly-reading}
 \E\tr R_r(A)\leq C(r)\tau_r(A)
 \quad\text{for all positive-definite }A.
\end{equation}
Indeed choose $r$ so that $2^r>C(r)$, then choose the finite tolerance
small enough to violate the proposed inequality. This conclusion is
independent of any ambiguity in an external statement.

Epperly's thesis \cite[Conjecture~11.2, printed p.~182]{Thesis} proposes
an improved polynomial $r$-step bound. The printed display is
\[
 \tr(A-A^{(r)})\leq C(r)\tr(A-[A]_r),\qquad C(r)=\operatorname{poly}(r).
\]
There are two notation issues: the thesis defines $A^{(r)}$ as a residual,
making the displayed subtraction an approximation trace, and the display
omits expectation. Its adjacent paragraph cites CETW25 Lemma~5.5 as
establishing the relevant bound with $C(r)=2^r$. That citation and the
surrounding discussion support the expected-residual interpretation
\eqref{eq:poly-reading}. \emph{Under that interpretation},
Theorem~\ref{thm:sharp} contradicts the proposed polynomial bound.
This repaired expression is an interpretation, not a verbatim quotation
of the printed equation, and the author's intended notation has not been
confirmed. The project source audit checked both the official deposit
and arXiv:2512.15929v1 of 17 December 2025. The primary claim of this note
remains sharpness of the independently stated published upper bound.

Every matrix constructed here has dimension $r+1$ and is exhausted by its
next pivot. The argument therefore does not refute the thesis's separate
oversampling conjecture, the $O(r/\varepsilon)$ question discussed in
\cite{NewAnalysis}, or valid oversampling guarantees. It also does not
establish a median-error or fixed-confidence obstruction. Rare paths can
matter to the expectation, and no practical application degradation is
claimed.

\paragraph{Attribution and status.}
The project was directed by Jitendra Prajapati. OpenAI Codex produced the
construction, derivation, explicit integer rescaling, implementation,
and source-audit work. Separately tasked AI agents checked the main
mechanism, uniform perturbation bounds, exponent arithmetic, integer
rescaling, finite-control artifacts, and source scope. These checks are
not formal verification or independent human mathematical validation.
Bounded literature searches did not establish priority; novelty and
publication priority remain unconfirmed. This public research draft has
not been peer reviewed and claims no external endorsement or journal
acceptance. Independent human mathematical review remains pending.

\begin{thebibliography}{9}
\bibitem{CETW}
Y.~Chen, E.~N.~Epperly, J.~A.~Tropp, and R.~J.~Webber.
\newblock Randomly pivoted Cholesky: Practical approximation of a kernel
matrix with few entry evaluations.
\newblock \emph{Communications on Pure and Applied Mathematics}
78(5) (2025), 995--1041.
\newblock \href{https://doi.org/10.1002/cpa.22234}{doi:10.1002/cpa.22234}.
Lemma~5.5 and Section~5.4.

\bibitem{Thesis}
E.~N.~Epperly.
\newblock \emph{Make the Most of What You Have: Resource-Efficient
Randomized Algorithms for Matrix Computations}.
\newblock PhD thesis, California Institute of Technology, 2025.
\newblock \href{https://doi.org/10.7907/pef3-mg80}{doi:10.7907/pef3-mg80}.
Chapter~11, Conjecture~11.2, printed p.~182.
Lightly edited revision:
\href{https://arxiv.org/abs/2512.15929v1}{arXiv:2512.15929v1},
17 December 2025.

\bibitem{DRVW}
A.~Deshpande, L.~Rademacher, S.~Vempala, and G.~Wang.
\newblock Matrix Approximation and Projective Clustering via Volume Sampling.
\newblock \emph{Theory of Computing} 2 (2006), 225--247.
\newblock \href{https://doi.org/10.4086/toc.2006.v002a012}
{doi:10.4086/toc.2006.v002a012}. Proposition~3.3.

\bibitem{DKM}
M.~Derezi\'{n}ski, R.~Khanna, and M.~W.~Mahoney.
\newblock Improved guarantees and a multiple-descent curve for Column Subset
Selection and the Nystr\"om method.
\newblock \href{https://arxiv.org/abs/2002.09073v3}{arXiv:2002.09073v3},
18 December 2020. Section~3.

\bibitem{Simons}
N.~Amsel et al.
\newblock Linear Systems and Eigenvalue Problems: Open Questions from a
Simons Workshop.
\newblock \href{https://arxiv.org/abs/2602.05394v2}{arXiv:2602.05394v2}, 2026.
Section~4.2 (scribe: Mark Fornace), equations~(16)--(17).

\bibitem{NewAnalysis}
E.~Epperly.
\newblock A new analysis of the randomly pivoted Cholesky algorithm.
\newblock \href{https://arxiv.org/abs/2608.20633}{arXiv:2608.20633v1}, 2026.
\end{thebibliography}
\end{document}
